\documentclass[12pt,oneside]{abntex2}

\usepackage[utf8]{inputenc}
\usepackage[T1]{fontenc}
\usepackage[scaled=1]{helvet}
\usepackage{indentfirst}
\usepackage{setspace}
\usepackage{pdflscape}
\usepackage[backend=biber,style=abnt]{biblatex}
\usepackage{graphicx}
\usepackage[table]{xcolor}
\addbibresource{referencias.bib}
\hypersetup{hidelinks}
\usepackage{float}
\renewcommand{\familydefault}{\sfdefault}

\OnehalfSpacing
\setlength{\parskip}{0.5\baselineskip}
\setsecheadstyle{\fontsize{14pt}{16pt}\selectfont\bfseries}
\setsubsecheadstyle{\fontsize{14pt}{16pt}\selectfont\bfseries}
\renewcommand{\chapnamefont}{\fontsize{14pt}{16pt}\selectfont\bfseries}
\renewcommand{\chapnumfont}{\fontsize{14pt}{16pt}\selectfont\bfseries}
\renewcommand{\chaptitlefont}{\fontsize{14pt}{16pt}\selectfont\bfseries}

\titulo{Desenvolvimento de um Laboratório de Detecção de Intrusão: Comparando Machine Learning e Deep Learning na Identificação de Ameaças em Redes}
\autor{EDUARDO DE MATTOS BORTOTO}
\local{São Paulo}
\data{2026}

\instituicao{%
  FACULDADE DE TECNOLOGIA DE SÃO PAULO - FATEC-SP
}

\tipotrabalho{Trabalho de conclusão de curso}
\orientador{Profª. Me. Grace Anne Pontes Borges}
\preambulo{Trabalho de conclusão do curso apresentado à FATEC-SP, como exigência parcial para a obtenção do grau de Tecnólogo em Análise e Desenvolvimento de Sistemas, sob orientação da Professora Mestre Grace Anne Pontes Borges.}

\makepagestyle{tccright}
\makeoddhead{tccright}{}{}{}
\makeevenhead{tccright}{}{}{}
\makeoddfoot{tccright}{}{}{\thepage}
\makeevenfoot{tccright}{}{}{\thepage}

\begin{document}

\pagenumbering{arabic}
\pagestyle{tccright}

\pretextual
\begin{capa}
\begin{center}
\bfseries\imprimirinstituicao

\vfill
\textbf{\imprimirautor}

\vfill
{\fontsize{16pt}{19pt}\selectfont\bfseries\imprimirtitulo\par}

\vfill
\imprimirlocal
\par
\imprimirdata
\end{center}
\end{capa}

\begin{folhaderosto}
\thispagestyle{empty}
\setcounter{page}{2}
\begin{center}
\textbf{\imprimirinstituicao}

\vfill
\textbf{\imprimirautor}

\vfill
{\fontsize{16pt}{19pt}\selectfont\bfseries\imprimirtitulo\par}

\vfill
\hfill
\begin{minipage}{0.5\textwidth}
\SingleSpacing
\imprimirpreambulo
\end{minipage}

\vfill
\imprimirlocal
\par
\imprimirdata
\end{center}
\end{folhaderosto}

\pdfbookmark[0]{\contentsname}{toc}
\tableofcontents*
\cleardoublepage

\textual
\OnehalfSpacing
\linespread{1.5}\selectfont
\pagestyle{tccright}
\aliaspagestyle{chapter}{tccright}

\chapter{INTRODUÇÃO}

A elaboração de estratégias de defesa cibernética robustas tornou-se uma exigência crítica e inadiável no cenário tecnológico contemporâneo. A digitalização acelerada dos processos corporativos, governamentais e sociais, catalisada pela adoção massiva de computação em nuvem e dispositivos móveis, trouxe consigo uma complexidade sem precedentes e vulnerabilidades sistêmicas que, quando exploradas, geram impactos econômicos, operacionais e reputacionais severos. Nesse contexto de hiperconectividade, a superfície de ataque expande-se continuamente, tornando a preservação da confidencialidade, integridade e disponibilidade da informação um desafio dinâmico que exige evolução constante frente a adversários cada vez mais organizados.

A análise dos relatórios recentes revela uma mudança de tendência nos prejuízos financeiros decorrentes de ciberataques. Estes dados indicam que as defesas tradicionais, baseadas puramente em regras estáticas e assinaturas, já não são suficientes para conter a sofisticação das ameaças modernas, resultando em custos globais de violação de dados que atingiram patamares históricos (IBM, 2024, p. 4).

Segundo o relatório Cost of a Data Breach Report 2025, o custo médio global de uma violação de dados apresentou uma queda para US\$4,44 milhões, representando uma redução de 9\% em relação ao ano anterior (IBM, 2025, p. 3). Essa diminuição é atribuída diretamente à redução do ciclo de vida das violações: o tempo médio para identificar e conter um incidente caiu para 241 dias (o menor patamar em nove anos), impulsionado fundamentalmente pelo uso extensivo de Inteligência Artificial e automação pelas equipes de segurança (IBM, 2025, p. 7).

No Brasil, o cenário acompanhou a tendência global de redução, com o custo médio de uma violação caindo para US\$1,22 milhão em 2025 (IBM, 2025, p. 6). Contudo, a complexidade dos ataques permanece crítica. O vetor de ataque inicial mais comum passou a ser o phishing, responsável por 16\% das violações, desbancando o uso de credenciais roubadas (IBM, 2025, p. 16). Além disso, ataques envolvendo comprometimento da cadeia de suprimentos (supply chain) tornaram-se o segundo vetor mais frequente e o que demanda maior tempo para identificação e contenção, com o tempo médio de 267 dias (IBM, 2025, p. 18).

Diante da falibilidade humana e da persistente escassez de profissionais qualificados, que ainda afeta 48\% das organizações (IBM, 2025, p. 44), a implementação de IDSs autônomos baseados em Inteligência Artificial (IA) consolida-se não apenas como uma vantagem, mas como uma necessidade estratégica. A capacidade de tecnologias como Machine Learning e Deep Learning de correlacionar eventos complexos e identificar anomalias sutis provou-se eficaz, gerando uma economia média de US\$1,9 milhão para organizações que utilizam essas tecnologias de forma extensiva em comparação àquelas que não as adotam (IBM, 2025, p. 4).

\section{HIPÓTESES}

A pesquisa parte da premissa de que algoritmos de Deep Learning, em virtude de sua arquitetura baseada em camadas ocultas e capacidade de extração automática de características (feature extraction), apresentarão taxas de detecção de intrusão superiores, com maior acurácia e menor incidência de falsos negativos, especialmente frente a ataques complexos e tráfego criptografado.

Em contrapartida, estima-se que o algoritmo Random Forest, pela sua natureza de conjunto de árvores de decisão, oferecerá um desempenho computacional mais eficiente, caracterizado por menores tempos de treinamento e classificação. Essa característica, aliada a uma maior interpretabilidade dos resultados, tende a torná-lo uma solução mais viável para ambientes com recursos de hardware limitados, ainda que possa apresentar uma ligeira perda de precisão na identificação de ameaças desconhecidas em comparação às redes neurais profundas.

\section{OBJETIVOS}

O objetivo deste trabalho consiste em desenvolver e implementar um laboratório experimental de detecção de intrusão para realizar uma análise comparativa de desempenho entre técnicas de Machine Learning (Random Forest) e Deep Learning (Redes Neurais Profundas), visando identificar qual abordagem oferece a melhor relação entre eficácia na detecção de ameaças e custo computacional em redes de computadores.

Para alcançar essa finalidade, o estudo buscará fundamentar teoricamente os conceitos de Sistemas de Detecção de Intrusão (IDS), abordando as diferenças entre detecção por assinatura e por anomalia. Na sequência, será preparado um ambiente de dados através do pré-processamento, limpeza e normalização de um dataset de tráfego de rede padronizado. A etapa prática envolverá a implementação, treinamento e validação dos modelos em cenários controlados, culminando na avaliação quantitativa dos resultados através de métricas como acurácia, precisão, revocação e análise da matriz de confusão.

\section{METODOLOGIA}

Este trabalho caracteriza-se como uma pesquisa aplicada, de natureza quantitativa e experimental. A metodologia proposta divide-se em quatro etapas sequenciais: Levantamento Bibliográfico, Preparação dos Dados, Implementação Experimental e Análise de Resultados.

\subsection{Levantamento Bibliográfico}

A primeira etapa consiste na revisão sistemática da literatura em bases de dados acadêmicas (Google Scholar, IEEE Xplore, Scopus), buscando trabalhos relacionados à aplicação de IA em segurança da informação, com foco em publicações dos últimos cinco anos. Serão consultadas normas técnicas e relatórios de mercado (como IBM Security e Cisco) para contextualizar o cenário de ameaças.

\subsection{Seleção e Pré-processamento do Dataset}

Será utilizado um dataset público de referência para intrusão de rede, preferencialmente o CICIDS2017 (disponibilizado pelo Canadian Institute for Cybersecurity), que contém tráfego benigno e os ataques mais comuns (Brute Force, DoS, Web Attack, Infiltration, Botnet, DDoS). Os dados passarão por um processo de tratamento (ETL), que inclui:

\begin{itemize}
\item Remoção de valores nulos ou infinitos.
\item Normalização dos dados numéricos (escala Min-Max) para garantir a convergência eficiente dos algoritmos, especialmente da Rede Neural.
\end{itemize}

\subsection{Implementação e Experimentação}

O laboratório virtual será desenvolvido utilizando a linguagem de programação Python. Para o modelo de Machine Learning, será utilizada a biblioteca Scikit-learn para implementar o algoritmo Random Forest. Para o modelo de Deep Learning, serão utilizadas as bibliotecas TensorFlow e Keras para a construção de uma Rede Neural Profunda (DNN) com múltiplas camadas densas. O ambiente de execução será configurado em notebooks (Jupyter ou Google Colab), permitindo a reprodutibilidade dos experimentos.

\subsection{Análise de Dados e Métricas}

A comparação entre os modelos não se baseará apenas na acurácia global, visto que datasets de segurança tendem a ser desbalanceados. A análise crítica focará na matriz de confusão, observando especialmente a taxa de falsos negativos (ataques que passaram despercebidos) e falsos positivos (alarmes falsos). Serão calculadas as métricas de precisão, recall e F1-score para cada classe de ataque, permitindo concluir qual modelo é mais robusto para tipos específicos de ameaças.

\chapter{SISTEMAS DE DETECÇÃO DE INTRUSÃO}

\section{Conceitos fundamentais de IDS}

[Apresentar definição de IDS, objetivos e contexto de uso em segurança de redes.]

\section{Abordagens de detecção}

[Comparar detecção por assinatura, por anomalia e modelo híbrido, destacando vantagens e limitações.]

\section{Desafios atuais de detecção em redes}

[Descrever desafios como alto volume de tráfego, falsos positivos e ataques inéditos.]

\chapter{ALGORITMO RANDOM FOREST}

\section{Fundamentos teóricos}

[Explicar árvores de decisão, bagging e votação por comitê no Random Forest.]

\section{Hiperparâmetros e comportamento do modelo}

[Descrever parâmetros principais (n\_estimators, max\_depth, etc.) e seus efeitos.] 

\section{Vantagens e limitações para IDS}

[Apontar interpretabilidade, robustez e limitações de generalização para classes não vistas.]

\chapter{REDES NEURAIS PROFUNDAS (DNN)}

\section{Fundamentos de Deep Learning}

[Introduzir perceptron multicamadas, camadas ocultas, funções de ativação e retropropagação.]

\section{Arquitetura DNN aplicada ao problema}

[Descrever escolhas de arquitetura, função de perda, otimizador e regularização.]

\section{Vantagens e limitações para IDS}

[Discutir capacidade de modelar padrões complexos e desafios de treinamento em dados desbalanceados.]

\chapter{TRABALHOS RELACIONADOS}

\section{Estudos com Random Forest em detecção de intrusão}

[Revisar trabalhos que utilizaram Random Forest em datasets como CICIDS2017.]

\section{Estudos com DNN em detecção de intrusão}

[Revisar trabalhos que utilizaram DNN para classificação de tráfego malicioso.]

\section{Comparações entre Machine Learning e Deep Learning}

[Apresentar resultados anteriores, lacunas e posicionamento deste TCC em relação à literatura.]

\chapter{BASE TEÓRICA PARA A MONTAGEM DO LABORATÓRIO}

\section{Escolha do dataset e critérios de seleção}

[Justificar a escolha do CICIDS2017 e descrever composição, classes e limitações.]

\section{Pipeline de preparação de dados}

[Descrever consolidacão dos arquivos, limpeza, tratamento de nulos/infinitos e normalização.]

\section{Definição dos cenários experimentais}

[Definir cenário completo, sem PortScan e sem XSS, com justificativa metodológica.]

\section{Métricas e critérios de avaliação}

[Justificar uso de acurácia, precisão, revocação, F1-score, médias macro/ponderada e matriz de confusão.]

\chapter{RESULTADOS DO LABORATÓRIO}

Este capítulo apresenta os resultados obtidos na etapa prática do laboratório de detecção de intrusão, com foco na comparação entre os modelos DNN e Random Forest em três cenários experimentais. 

\section{Execução do Laboratório Experimental}

O laboratório foi estruturado em duas fases principais: tratamento dos dados e treinamento/avaliação dos modelos. Na primeira fase, os arquivos do CICIDS2017 foram consolidados em apenas 1 CSV, tratados e normalizados. Na segunda fase, os modelos foram treinados e avaliados em três cenários:

\begin{itemize}
\item Cenário completo: todas as classes presentes no dataset de teste, estavam presentes também no dataset de treinamento.
\item Cenário sem PortScan no treino: classe PortScan foi removida do dataset de treinamento e mantida apenas no dataset de teste para simular um ataque não visto durante o treinamento.
\item Cenário sem XSS no treino: classe XSS removida do dataset de treinamento e mantida no apenas dataset de teste para novamente simular um ataque não visto durante o treinamento.
\end{itemize}

No pré-processamento, o dataset unificado atingiu dimensão de 2.830.743 registros e 79 colunas. Após a divisão principal, foram gerados conjuntos com 1.979.513 amostras para treino e 848.363 para teste no cenário completo. No cenário sem PortScan, 110.968 registros dessa classe foram movidos do treino para o teste, resultando em 1.868.545 amostras para treino e 959.331 para teste. Já no cenário sem XSS, a divisão manteve-se similar ao cenário completo, com 1.979.513 amostras para treino e 848.822 para teste.

\section{Resultados dos Modelos no Cenário Completo}

No cenário completo, os dois modelos apresentaram desempenho elevado. A DNN alcançou acurácia global de 0,98, enquanto o Random Forest atingiu acurácia global de 1,00. Esses resultados indicam que, quando todas as classes estão representadas no treinamento, ambos os modelos são capazes de aprender padrões discriminativos de forma consistente.

Apesar da vantagem numérica do Random Forest em acurácia global, essa métrica isolada não resume todo o comportamento por classe. Por esse motivo, a análise foi orientada também por matriz de confusão e métricas por classe (precisão, recall e F1-score), permitindo observar possíveis confusões entre tipos de ataque específicos.

Além das métricas de classificação, foi realizada a comparação do tempo de inferência no dataset de teste como proxy de custo computacional em produção. No cenário completo (848.363 amostras), a DNN levou 17,3258 segundos para classificar o conjunto de teste, enquanto o Random Forest levou 3,3474 segundos. Em termos práticos, o Random Forest foi aproximadamente 5,18 vezes mais rápido nesse cenário.

\begin{table}[H]
\centering
\small
\label{tab:comparacao_completo}
\begin{tabular}{lcc}
  \toprule
  \textbf{Modelo} & \textbf{Acurácia} & \textbf{Tempo de classificação (s)} \\
  \midrule
  DNN & 0,98 & 17,3258 \\
  Random Forest & 1,00 & 3,3474 \\
  \bottomrule
\end{tabular}
\caption{Comparação de desempenho e tempo de inferência no cenário completo}
\end{table}

Para ampliar a comparação do cenário completo, também foram consolidadas as métricas de precisão, recall e F1-score com base no relatório de classificação de cada modelo. A Tabela \ref{tab:metricas_completo} apresenta os resultados para as médias macro e ponderada (weighted), permitindo comparar o equilíbrio de desempenho entre classes e o desempenho global ponderado pelo suporte.

\begin{table}[H]
\centering
\small
\label{tab:metricas_completo}
\begin{tabular}{llccc}
  \toprule
  \textbf{Modelo} & \textbf{Métrica} & \textbf{Precisão} & \textbf{Revocação} & \textbf{F1-score} \\
  \midrule
  DNN & Acurácia & - & 0.98 & - \\
  DNN & Média Macro & 0.83 & 0.64 & 0.67 \\
  DNN & Média Ponderada & 0.98 & 0.98 & 0.98 \\
  \midrule
  Random Forest & Acurácia & - & 1.00 & - \\
  Random Forest & Média Macro & 0.94 & 0.83 & 0.86 \\
  Random Forest & Média Ponderada & 1.00 & 1.00 & 1.00 \\
  \bottomrule
\end{tabular}
\caption{Métricas agregadas no cenário completo}
\end{table}

Em relação à matriz de confusão do cenário completo, os dois modelos apresentaram concentração de acertos na diagonal principal. Para tornar a análise objetiva no corpo do texto, a Tabela \ref{table:mc_dnn_completa} apresenta a matriz de confusão completa da DNN (15x15), e a Tabela \ref{table:mc_rf_completa} do Random Forest.

\newcommand{\ok}[1]{\cellcolor{green!25}#1}
\newcommand{\err}[1]{\cellcolor{red!25}#1}

% \begin{landscape}
\begin{table}[H]
\centering
\scriptsize
\resizebox{\linewidth}{!}{
\begin{tabular}{l|rrrrrrrrrrrrrrr}
\hline
 & BENIGN & Bot & DDoS & DoS GE & Hulk & Slowhttp & Slowloris & FTP & Heart & Infil & PortScan & SSH & Brute & SQL & XSS \\
\hline
Real\_BENIGN & \ok{671956} & \err{13} & \err{15} & \err{25} & \err{3737} & \err{30} & \err{12} & \err{1} & 0 & 0 & \err{5197} & \err{27} & 0 & 0 & \err{1} \\
Real\_Bot & \err{354} & \ok{210} & 0 & 0 & 0 & 0 & 0 & 0 & 0 & 0 & 0 & 0 & 0 & 0 & 0 \\
Real\_DDoS & \err{527} & 0 & \ok{37776} & \err{3} & \err{1} & 0 & 0 & 0 & 0 & 0 & 0 & 0 & 0 & 0 & 0 \\
Real\_DoS GE & \err{85} & 0 & 0 & \ok{3074} & 0 & \err{2} & \err{2} & 0 & 0 & 0 & 0 & 0 & 0 & 0 & 0 \\
Real\_DoS Hulk & \err{76} & 0 & 0 & \err{1} & \ok{69190} & 0 & 0 & 0 & 0 & 0 & 0 & 0 & 0 & 0 & 0 \\
Real\_Slowhttp & \err{17} & 0 & 0 & 0 & 0 & \ok{1626} & \err{14} & 0 & 0 & 0 & 0 & \err{1} & 0 & 0 & 0 \\
Real\_Slowloris & \err{11} & 0 & 0 & \err{1} & \err{1} & \err{37} & \ok{1704} & \err{2} & 0 & 0 & 0 & \err{1} & 0 & 0 & 0 \\
Real\_FTP & 0 & 0 & 0 & 0 & \err{2} & 0 & \err{4} & \ok{2320} & 0 & 0 & 0 & 0 & 0 & 0 & 0 \\
Real\_Heart & \err{3} & 0 & 0 & 0 & 0 & 0 & 0 & 0 & \ok{3} & 0 & 0 & 0 & 0 & 0 & 0 \\
Real\_Infil & \err{10} & 0 & 0 & 0 & 0 & 0 & 0 & 0 & 0 & \ok{0} & 0 & 0 & 0 & 0 & 0 \\
Real\_PortScan & \err{2970} & 0 & 0 & 0 & \err{13} & 0 & 0 & 0 & 0 & 0 & \ok{44850} & \err{3} & 0 & 0 & 0 \\
Real\_SSH & \err{200} & 0 & 0 & 0 & \err{2} & 0 & 0 & \err{4} & 0 & 0 & 0 & \ok{1615} & 0 & 0 & 0 \\
Real\_Brute & \err{388} & 0 & 0 & 0 & 0 & 0 & 0 & 0 & 0 & 0 & 0 & \err{34} & \ok{13} & 0 & 0 \\
Real\_SQL & \err{3} & 0 & 0 & \err{2} & 0 & 0 & 0 & 0 & 0 & 0 & 0 & \err{1} & 0 & \ok{0} & 0 \\
Real\_XSS & \err{178} & 0 & 0 & \err{2} & \err{7} & 0 & 0 & 0 & 0 & 0 & 0 & \err{1} & \err{1} & 0 & \ok{4} \\
\hline
\end{tabular}
}
\caption{Matriz de Confusão - DNN}
\label{table:mc_dnn_completa}
\end{table}
% \end{landscape}

\begin{table}[H]
\centering
\small
\begin{tabular}{lrrrr}
\toprule
Classe & Precisão & Revocação & F1-score & Suporte \\
\midrule
BENIGN & 0.99 & 0.99 & 0.99 & 681014 \\
Bot & 0.94 & 0.37 & 0.53 & 564 \\
DDoS & 1.00 & 0.99 & 0.99 & 38307 \\
DoS GoldenEye & 0.99 & 0.97 & 0.98 & 3163 \\
DoS Hulk & 0.95 & 1.00 & 0.97 & 69267 \\
DoS Slowhttptest & 0.96 & 0.98 & 0.97 & 1658 \\
DoS slowloris & 0.98 & 0.97 & 0.98 & 1757 \\
FTP-Patator & 1.00 & 1.00 & 1.00 & 2326 \\
Heartbleed & 1.00 & 0.50 & 0.67 & 6 \\
Infiltration & 0.00 & 0.00 & 0.00 & 10 \\
PortScan & 0.90 & 0.94 & 0.92 & 47836 \\
SSH-Patator & 0.96 & 0.89 & 0.92 & 1821 \\
Web Attack - Brute Force & 0.93 & 0.03 & 0.06 & 435 \\
Web Attack - Sql Injection & 0.00 & 0.00 & 0.00 & 6 \\
Web Attack - XSS & 0.80 & 0.02 & 0.04 & 193 \\
\midrule
\textbf{Acurácia} & - & 0.98 & -  & 848363 \\
\textbf{Média Macro} & 0.83 & 0.64 & 0.67 & 848363 \\
\textbf{Média Ponderada} & 0.98 & 0.98 & 0.98 & 848363 \\
\bottomrule
\end{tabular}
\caption{Relatório de Métricas do Modelo DNN}
\label{table:rel_metr_DNN}
\end{table}

% \begin{landscape}
\begin{table}[H]
\centering
\scriptsize
\resizebox{\linewidth}{!}{
\begin{tabular}{l|rrrrrrrrrrrrrrr}
\hline
 & BENIGN & Bot & DDoS & DoS GE & Hulk & Slowhttp & Slowloris & FTP & Heart & Infil & PortScan & SSH & Brute & SQL & XSS \\
\hline
Real\_BENIGN & \ok{680465} & \err{49} & \err{5} & \err{1} & \err{193} & \err{11} & 0 & 0 & 0 & 0 & \err{288} & 0 & \err{2} & 0 & 0 \\
Real\_Bot & \err{285} & \ok{279} & 0 & 0 & 0 & 0 & 0 & 0 & 0 & 0 & 0 & 0 & 0 & 0 & 0 \\
Real\_DDoS & \err{18} & 0 & \ok{38289} & 0 & 0 & 0 & 0 & 0 & 0 & 0 & 0 & 0 & 0 & 0 & 0 \\
Real\_DoS GE & \err{5} & 0 & 0 & \ok{3151} & \err{4} & \err{3} & 0 & 0 & 0 & 0 & 0 & 0 & 0 & 0 & 0 \\
Real\_DoS Hulk & \err{184} & 0 & 0 & \err{6} & \ok{69073} & 0 & 0 & 0 & 0 & 0 & \err{4} & 0 & 0 & 0 & 0 \\
Real\_Slowhttp & \err{6} & 0 & 0 & 0 & 0 & \ok{1644} & \err{8} & 0 & 0 & 0 & 0 & 0 & 0 & 0 & 0 \\
Real\_Slowloris & \err{4} & 0 & 0 & 0 & 0 & \err{3} & \ok{1750} & 0 & 0 & 0 & 0 & 0 & 0 & 0 & 0 \\
Real\_FTP & 0 & 0 & 0 & 0 & 0 & 0 & 0 & \ok{2326} & 0 & 0 & 0 & 0 & 0 & 0 & 0 \\
Real\_Heart & \err{1} & 0 & 0 & 0 & 0 & 0 & 0 & 0 & \ok{5} & 0 & 0 & 0 & 0 & 0 & 0 \\
Real\_Infil & \err{2} & 0 & 0 & 0 & 0 & 0 & 0 & 0 & 0 & \ok{8} & 0 & 0 & 0 & 0 & 0 \\
Real\_PortScan & \err{9} & 0 & 0 & 0 & \err{8} & 0 & 0 & 0 & 0 & 0 & \ok{47818} & 0 & \err{1} & 0 & 0 \\
Real\_SSH & \err{2} & 0 & 0 & 0 & \err{1} & 0 & 0 & 0 & 0 & 0 & 0 & \ok{1818} & 0 & 0 & 0 \\
Real\_Brute & \err{9} & 0 & 0 & 0 & 0 & 0 & 0 & 0 & 0 & 0 & 0 & 0 & \ok{352} & 0 & \err{74} \\
Real\_SQL & \err{4} & 0 & 0 & \err{1} & 0 & 0 & 0 & 0 & 0 & 0 & 0 & 0 & 0 & \ok{1} & 0 \\
Real\_XSS & \err{13} & 0 & 0 & 0 & 0 & 0 & 0 & 0 & 0 & 0 & \err{1} & 0 & \err{107} & 0 & \ok{72} \\
\hline
\end{tabular}
}
\caption{Matriz de Confusão - Random Forest}
\label{table:mc_rf_completa}
\end{table}
% \end{landscape}

\begin{table}[H]
\centering
\small
\begin{tabular}{lrrrr}
\toprule
Classe & Precisão & Revocação & F1-score & Suporte \\
\midrule
BENIGN & 1.00 & 1.00 & 1.00 & 681014 \\
Bot & 0.85 & 0.49 & 0.63 & 564 \\
DDoS & 1.00 & 1.00 & 1.00 & 38307 \\
DoS GoldenEye & 1.00 & 1.00 & 1.00 & 3163 \\
DoS Hulk & 1.00 & 1.00 & 1.00 & 69267 \\
DoS Slowhttptest & 0.99 & 0.99 & 0.99 & 1658 \\
DoS slowloris & 1.00 & 1.00 & 1.00 & 1757 \\
FTP-Patator & 1.00 & 1.00 & 1.00 & 2326 \\
Heartbleed & 1.00 & 0.83 & 0.91 & 6 \\
Infiltration & 1.00 & 0.80 & 0.89 & 10 \\
PortScan & 0.99 & 1.00 & 1.00 & 47836 \\
SSH-Patator & 1.00 & 1.00 & 1.00 & 1821 \\
Web Attack - Brute Force & 0.76 & 0.81 & 0.78 & 435 \\
Web Attack - Sql Injection & 1.00 & 0.17 & 0.29 & 6 \\
Web Attack - XSS & 0.49 & 0.37 & 0.42 & 193 \\
\midrule
\textbf{Acurácia} & - & 1.00 & - & 848363 \\
\textbf{Média Macro} & 0.94 & 0.83 & 0.86 & 848363 \\
\textbf{Média Ponderada} & 1.00 & 1.00 & 1.00 & 848363 \\
\bottomrule
\end{tabular}
\caption{Relatório de Métricas do Modelo random Forest}
\label{table:rel_metr_rf}
\end{table}

As tabelas evidenciam que o Random Forest apresentou diagonal principal mais forte nas classes críticas do cenário completo. Em \textit{PortScan}, o Random Forest acertou 47.818 de 47.836 registros na diagonal (aprox. 99,96\%), enquanto a DNN acertou 44.850 de 47.836 (aprox. 93,75\%). Em \textit{Web Attack - Brute Force}, o Random Forest registrou 352 acertos de 435 (aprox. 80,92\%), ao passo que a DNN registrou 13 de 435 (aprox. 2,99\%). O mesmo padrão aparece em \textit{Web Attack - XSS}, com 72 acertos de 193 para Random Forest (aprox. 37,31\%) contra 4 de 193 para DNN (aprox. 2,07\%).

\section{Resultados em Cenários com Classe Ausente no Treinamento}

Nos cenários de generalização, os resultados foram mais heterogêneos e revelaram diferenças importantes entre os modelos.

\subsection{Análise do cenário sem PortScan no treinamento}

No experimento sem PortScan no treinamento, os modelos apresentaram queda de desempenho em relação ao cenário completo, indicando dificuldade de generalização quando esta classe é totalmente ausente no treino. A DNN obteve acurácia global de 0,79, enquanto o Random Forest alcançou 0,83, mantendo vantagem também neste cenário.

Os dados dos relatórios de classificação confirmam que o impacto não ficou restrito ao PortScan. Na DNN, a média macro caiu de 0,83/0,64/0,67 (precisão/recall/F1) no cenário completo para 0,52/0,42/0,45 sem PortScan, e a média ponderada caiu de 0,98/0,98/0,98 para 0,70/0,79/0,74. Em classes não omitidas, houve queda importante de recall, como em DoS Slowhttptest (0,98 para 0,25) e SSH-Patator (0,89 para 0,49), mostrando piora de generalização além da classe removida.

No Random Forest, a acurácia também caiu (1,00 para 0,83), com redução das médias de 0,94/0,83/0,86 para 0,86/0,77/0,79 (macro) e de 1,00/1,00/1,00 para 0,70/0,83/0,76 (ponderada). Mesmo mantendo bom desempenho em algumas classes, houve perda relevante em BENIGN (precisão de 1,00 para 0,81), indicando aumento de confusão após a retirada de PortScan do treino. Esses valores sustentam que a remoção de PortScan piora o resultado geral e também afeta classes correlatas.

\begin{center}
\textit{[Inserir tabela de resultados do cenário sem PortScan]} 
\end{center}

\subsection{Análise do cenário sem XSS no treinamento}

No experimento sem XSS no treinamento, o impacto global foi menor. A DNN manteve acurácia de 0,98 e o Random Forest permaneceu com acurácia de 1,00, indicando maior estabilidade quando a classe ausente possui menor representatividade e menor influência sobre o desempenho agregado.

Para verificar se houve o mesmo efeito do cenário sem PortScan (piora generalizada), a comparação foi feita em classes não relacionadas diretamente ao XSS. Na DNN, a média macro variou de 0,83/0,64/0,67 para 0,78/0,64/0,67, enquanto a média ponderada permaneceu em torno de 0,99/0,98/0,98. Em classes relevantes não omitidas, as variações foram pequenas: PortScan teve recall de 0,94 para 0,93 e DoS Slowhttptest manteve recall de 0,98 para 0,98, sugerindo ausência de degradação ampla entre classes.

No Random Forest, apesar da acurácia permanecer em 1,00 e da média ponderada ficar em 1,00/1,00/1,00, a média macro caiu de 0,94/0,83/0,86 para 0,88/0,82/0,82. Ainda assim, o comportamento das classes majoritárias permaneceu estável (por exemplo, PortScan com recall de 1,00 para 1,00 e BENIGN com 1,00 para 1,00), enquanto a deterioração concentrou-se em classes web: Web Attack - Brute Force teve precisão reduzida de 0,76 para 0,42 (F1 de 0,78 para 0,58), e Web Attack - XSS passou para 0,00/0,00/0,00 no cenário sem treino da própria classe.

Esse comportamento é compatível com o desbalanceamento do dataset. No cenário sem XSS, a classe Web Attack - XSS representa apenas 652 amostras em 848.822 no teste (aprox. 0,08\%), enquanto BENIGN sozinho concentra 681.014 amostras. Com essa distribuição, a acurácia global e a média ponderada ficam dominadas pelas classes majoritárias; por isso, mesmo com desempenho nulo em XSS (recall 0,00 no Random Forest), o valor global pode permanecer em 1,00.

Assim, a resposta para este cenário é: não houve o mesmo padrão de piora generalizada observado na retirada de PortScan. A remoção de XSS provocou piora localizada, principalmente no próprio XSS e em classes web correlatas, sem colapso amplo da detecção nas demais classes.

\begin{center}
\textit{[Inserir tabela de resultados do cenário sem XSS]}
\end{center}

Comparando os dois casos, conclui-se que a retirada de classes do treinamento pode sim piorar o resultado no geral, porém com intensidade diferente conforme a classe omitida. A remoção de PortScan gerou impacto mais amplo e mais forte no desempenho global e entre classes, enquanto a remoção de XSS produziu efeito predominantemente pontual.

\section{Comparação Geral de Tempo de Classificação}

De forma complementar aos resultados de acurácia, foi consolidada a comparação de tempo de classificação entre os três cenários testados. Essa análise permite estimar o custo computacional de inferência dos modelos em contextos distintos.

O padrão observado foi consistente: o Random Forest manteve tempos de inferência menores em todos os casos, enquanto a DNN apresentou maior tempo de classificação, principalmente no cenário sem PortScan no treinamento.

\begin{table}[htb]
\centering
\caption{Tempo de classificação por cenário (proxy de custo computacional)}
\label{tab:tempos_cenarios}
\begin{tabular}{|l|c|c|c|}
\hline
\textbf{Cenário} & \textbf{DNN (s)} & \textbf{Random Forest (s)} & \textbf{Amostras no teste} \\
\hline
Completo & 17,3258 & 3,3474 & 848.363 \\
\hline
Sem PortScan no treino & 19,5003 & 3,4952 & 959.331 \\
\hline
Sem XSS no treino & 16,8805 & 3,0954 & 848.822 \\
\hline
\end{tabular}
\end{table}

\section{Síntese dos Achados do Laboratório}

Os resultados apontam que os dois modelos são eficazes em cenário fechado, porém respondem de maneira distinta quando há ausência de classes no treinamento. O cenário sem PortScan foi o mais desafiador e destacou limitações de generalização, especialmente para a DNN. Já os cenários com XSS mantiveram desempenho global elevado, indicando menor impacto dessa remoção na configuração testada.

Embora a DNN seja frequentemente descrita como mais capaz de generalizar para padrões complexos e ataques não vistos, esse comportamento não foi observado de forma consistente nos experimentos deste trabalho. Nos cenários com classe ausente no treinamento, o Random Forest manteve desempenho mais estável e apresentou melhores resultados agregados.

\postextual
\nocite{*}
\printbibliography[title={REFERÊNCIAS BIBLIOGRÁFICAS}]

\end{document}